\chapter{Stratégies crypto à moyenne fréquence}\label{ch:s1:crypto-medium-frequency-strategies}

En 2021, le financement du contrat perpétuel Bitcoin de Binance s'élevait en moyenne à 30.6~\% par an, versé par les positions longues aux positions courtes~: un book long au comptant sur le bitcoin et short sur le contrat perpétuel,
sans opinion sur le prix, l'encaissait. En novembre 2022, FTX.com, l'une des plus grandes plateformes sur lesquelles de tels books étaient gérés, a cessé de honorer les retraits~;
à la requête en faillite, ses débiteurs ont localisé des actifs valant environ un cinquième de ce que les clients leur devaient. Les stratégies crypto à moyenne fréquence sont fondées sur le carry, le momentum
et les flux, comme sur les autres marchés, avec un ajout~: la plateforme elle-même peut faire défaut. Sur les données de Binance, le book carry a rapporté un rendement de 25.3~\% sur le capital en 2021 et
3.3~\% en 2022~; un book réparti sur quatre plateformes, chacune faisant défaut avec une probabilité de 5~\% par an et perdant ce qu'a perdu FTX, abandonne 20.2~\% de son capital lorsque
l'une d'elles fait défaut. La construction correspondante est \texttt{firm.cryptomf}.

\section{Carry de financement et carry de base}

\begin{definition}[Carry de financement, carry de base]\label{def:s1:crypto-medium-frequency-strategies:carry}
Le \emph{carry de financement}\index{carry de financement} est le rendement d'une position longue au comptant et courte sur un contrat perpétuel portant sur la même cryptomonnaie, qui perçoit les
paiements de financement du contrat lorsqu'ils sont positifs (les longs paient les shorts) et les verse lorsqu'ils sont négatifs, sans exposition nette au prix. Le \emph{carry de
base}\index{carry de base} est la même opération avec un contrat à échéance~: position longue au comptant et courte sur le contrat, qui perçoit la prime du contrat sur le comptant à mesure qu'il converge
à l'échéance.
\end{definition}

\begin{dated}{2026-09}{Financement du contrat perpétuel sur Binance}\label{dat:s1:crypto-medium-frequency-strategies:funding}
Le taux de financement de Binance pour un contrat perpétuel est la prime moyenne du contrat sur son indice, augmentée d'un terme d'intérêt borné, versé toutes les huit heures à
00:00, 08:00 et 16:00 UTC~; pour la plupart des contrats, le terme d'intérêt vaut 0.01~\% par intervalle. Le livre~3, chapitre~17 décrit le mécanisme et son
indice de prime.
\end{dated}

Le rendement sur capital du book carry (le \cref{lst:s1:crypto-medium-frequency-strategies:data}) est le financement annuel, diminué des frais de transaction (0.2~\% par an,
hypothèse), divisé par le capital immobilisé~: la position au comptant plus la marge du contrat perpétuel (20~\% du notionnel, hypothèse).

\begin{center}
\small
\begin{tabular}{@{}lccccccc@{}}
\toprule
année & 2020 & 2021 & 2022 & 2023 & 2024 & 2025 & 2026\\
\midrule
financement, moyenne annualisée & 17.19~\% & 30.61~\% & 4.16~\% & 7.87~\% & 11.92~\% & 5.13~\% & 2.94~\%\\
part des taux négatifs & 14.3~\% & 7.3~\% & 22.1~\% & 10.1~\% & 8.4~\% & 12.9~\% & 26.1~\%\\
book carry, rendement sur capital & 14.2~\% & 25.3~\% & 3.3~\% & 6.4~\% & 9.8~\% & 4.1~\% & 2.3~\%\\
\bottomrule
\end{tabular}
\end{center}

2026 court jusqu'au 24 septembre. Le carry suit la demande de levier du marché (la \cref{fig:s1:crypto-medium-frequency-strategies:funding})~: élevé sur le
marché haussier de 2021, faible après le krach de 2022, en reprise en 2024. Schmeling, Schrimpf et Todorov observent le même phénomène dans la base des contrats à échéance, parfois
au-dessus de 40~\% par an, et l'attribuent à la demande de levier d'investisseurs plus petits et suiveurs de tendance, se heurtant au capital d'arbitrage limité par des
frictions réglementaires et de marge. Le carry est le prix de cette rareté. En 2024, il était proche du terme d'intérêt de 0.01~\% par intervalle ($0.0001 \times 3 \times 365 = 11~\%$ par
an)~; depuis lors, il est resté en dessous.

\begin{omfigure}
\begin{tikzpicture}
\begin{axis}[width=10.5cm,height=5cm,ybar,bar width=10pt,xlabel={\small année},ylabel={\small \% par an},tick label style={font=\footnotesize},
  symbolic x coords={2020,2021,2022,2023,2024,2025,2026},xtick=data,ymin=0,ymax=35,enlarge x limits=0.08,grid=major,grid style={gray!25},
  legend columns=2,legend style={font=\scriptsize,at={(0.5,-0.3)},anchor=north,draw=none,fill=none,
  /tikz/every even column/.append style={column sep=8pt}},area legend,table/col sep=comma]
\addplot[fill=omExo!50,draw=omExo] table[x=year,y=funding_pct]{figdata/strategies-1/27-crypto-medium-frequency-strategies/funding.csv};
\addplot[fill=omThm!60,draw=omThm] table[x=year,y=carry_pct]{figdata/strategies-1/27-crypto-medium-frequency-strategies/funding.csv};
\legend{financement moyen (annualisé),rendement sur capital du book carry}
\end{axis}
\end{tikzpicture}

\omcaption{Contrat perpétuel BTCUSDT de Binance~: le taux de financement moyen par année civile, annualisé, et le rendement sur capital d'un book long au comptant et short sur le
contrat perpétuel (marge de 20~\%, coûts de 0.2~\% par an, hypothèses)~; 2026 jusqu'au 24 septembre. Données~: point d'accès du taux de financement de Binance, via le livre~3~;
\texttt{s1\_crypto.funding}.}\label{fig:s1:crypto-medium-frequency-strategies:funding}
\end{omfigure}

\section{Momentum transversal sur les cryptomonnaies}

Liu et Tsyvinski ont établi que les rendements des cryptomonnaies ne sont pas exposés aux facteurs boursiers et macroéconomiques habituels, mais sont prédits par des facteurs
propres à la crypto~: un fort effet de \omterm{def:s1:trend-following:trend}{momentum temporel} (time-series momentum) et des indicateurs de substitution de l'attention des investisseurs. Liu, Tsyvinski et Wu ont établi que trois facteurs, le marché crypto,
la taille et le momentum, expliquent la dimension transversale des rendements attendus des cryptomonnaies. \texttt{firm.cryptomf} simule six années de 50 cryptomonnaies avec un facteur de marché de 60~\%
de marché, une volatilité spécifique de 80~\% avec queues épaisses (une loi de Student-$t$ à trois degrés de liberté) et une dérive persistante injectée, d'une demi-vie de 60
jours, et fait tourner un book long--short hebdomadaire sur les rendements passés, quintile supérieur contre quintile inférieur~:

\begin{center}
\small
\begin{tabular}{@{}lccc@{}}
\toprule
fenêtre historique & 7 jours & 28 jours & 56 jours\\
\midrule
ratio de Sharpe avant coûts & 0.85 & 1.56 & 1.29\\
après 10 pb par unité échangée & 0.42 & 1.35 & 1.13\\
après 20 pb par unité échangée & $-0.01$ & 1.13 & 0.97\\
\bottomrule
\end{tabular}
\end{center}

Une semaine d'historique n'est surtout que du bruit à ces volatilités~; quatre semaines mesurent mieux la dérive et échangent moins. Les coûts en crypto varient fortement selon la plateforme et
la cryptomonnaie, et les petites cryptomonnaies où le momentum est le plus fort sont les plus coûteuses à échanger.

\section{Signaux on-chain}

\begin{definition}[Signal on-chain]\label{def:s1:crypto-medium-frequency-strategies:onchain}
Un \emph{signal on-chain}\index{signal on-chain} est un signal de trading construit à partir des registres publics d'une blockchain~: flux de cryptomonnaies vers et depuis les
adresses connues des plateformes, soldes des gros détenteurs, activité des adresses, émissions et destructions de jetons, accessibles à quiconque lit la chaîne (livre~3,
chapitre~26).
\end{definition}

L'hypothèse classique~: les cryptomonnaies envoyées vers les plateformes sont sur le point d'être vendues. Le chapitre l'injecte~: l'entrée nette de chaque cryptomonnaie, standardisée, présente une corrélation de
$-0.03$ avec le rendement spécifique du lendemain. Mesurée comme IC de rang sur les 50 cryptomonnaies, l'entrée centrée réduite (par rapport aux 30 derniers jours) donne $-0.031$
($t = -10.2$) sur six ans. L'effet injecté est retrouvé. Dans les données réelles, la correspondance entre adresses et plateformes est incomplète et évolue, et les flux
entre les \omterm{def:s1:residual-and-principal-component-stat-arb:eigenportfolio}{portefeuilles propres} à une plateforme ressemblent à des dépôts~; l'ossature de \texttt{firm.cryptomf} est l'endroit où un tel signal est testé, et non la preuve
qu'un tel signal fonctionne.

\section{Risque de plateforme et risque de conservation}

\begin{definition}[Risque de plateforme]\label{def:s1:crypto-medium-frequency-strategies:venue}
Le \emph{risque de plateforme}\index{risque de plateforme} est le risque de perdre des actifs détenus sur une plateforme de négociation ou chez un conservateur, à la suite de sa défaillance, d'une fraude, d'un piratage ou du gel des
retraits~; dans la crypto, il est supporté par le client, car les plateformes détiennent souvent elles-mêmes les actifs de leurs clients.
\end{definition}

L'analyse menée par les débiteurs de FTX.com donne l'ordre de grandeur d'une telle perte. À la requête en faillite, les dettes envers les clients s'élevaient à environ 11.2~milliards de dollars et les actifs localisés à
environ 2.2~milliards de dollars~: un déficit de 80.8~\%. Pendant les onze jours de novembre 2022 précédant la requête, les clients ont retiré 7.0~milliards de dollars de plus que
ce qu'ils avaient déposé, dont 3.3~milliards de dollars nets le 7 novembre à eux seuls. Le livre~3, chapitre~15 raconte l'épisode. Pour un book carry, la leçon est purement arithmétique
(le \cref{lst:s1:crypto-medium-frequency-strategies:venues})~: avec des plateformes qui font défaut de façon indépendante, avec une probabilité $p$ par an, et qui perdent une part $L$ des
actifs qui y sont détenus, la perte espérée est $pL$ quel que soit le nombre de plateformes, et la répartition des actifs ne modifie que sa distribution.

\begin{center}
\small
\begin{tabular}{@{}lcccc@{}}
\toprule
plateformes (parts égales), $p = 5~\%$, $L = 80.8~\%$ & 1 & 2 & 4 & 8\\
\midrule
perte lorsqu'une plateforme fait défaut (part du capital) & 80.8~\% & 40.4~\% & 20.2~\% & 10.1~\%\\
probabilité d'au moins une défaillance dans l'année & 5.0~\% & 9.9~\% & 18.7~\% & 34.0~\%\\
perte espérée annuelle & 4.1~\% & 4.1~\% & 4.1~\% & 4.1~\%\\
centile 99 de la perte annuelle & 80.8~\% & 40.4~\% & 40.4~\% & 20.2~\%\\
\bottomrule
\end{tabular}
\end{center}

Face au carry, une perte espérée d'environ 4~\% par an est déterminante~: les 25.3~\% de 2021 deviendraient un rendement net de 21.3~\%, et les 3.3~\% de 2022 seraient une perte
(la \cref{fig:s1:crypto-medium-frequency-strategies:venues}). La diversification entre plateformes rend supportable la mauvaise année~; elle ne rend pas la stratégie meilleure.
Ce qui réduit la perte espérée, c'est le choix des plateformes et de la conservation des actifs~: conservation séparée, actifs détenus hors de la bourse tant qu'ils ne sont pas nécessaires, et la
rapidité avec laquelle on peut retirer ses avoirs lorsqu'une plateforme se fragilise.

\begin{omfigure}
\begin{tikzpicture}
\begin{axis}[width=10cm,height=5cm,ybar,bar width=12pt,xlabel={\small nombre de plateformes},ylabel={\small~\% du capital, ou probabilité},
  tick label style={font=\footnotesize},symbolic x coords={1,2,4,8},xtick=data,ymin=0,ymax=90,enlarge x limits=0.2,grid=major,
  grid style={gray!25},legend columns=3,legend style={font=\scriptsize,at={(0.5,-0.3)},anchor=north,draw=none,fill=none,
  /tikz/every even column/.append style={column sep=8pt}},area legend,table/col sep=comma]
\addplot[fill=omThm!60,draw=omThm] table[x=k,y=one_pct]{figdata/strategies-1/27-crypto-medium-frequency-strategies/venues.csv};
\addplot[fill=omDef!60,draw=omDef] table[x=k,y=any_pct]{figdata/strategies-1/27-crypto-medium-frequency-strategies/venues.csv};
\addplot[fill=omExo!50,draw=omExo] table[x=k,y=mean_pct]{figdata/strategies-1/27-crypto-medium-frequency-strategies/venues.csv};
\legend{perte lorsqu'une plateforme fait défaut,probabilité d'une défaillance dans l'année,perte espérée annuelle}
\end{axis}
\end{tikzpicture}

\omcaption{Le capital d'un book carry est réparti à parts égales entre une et huit plateformes, chacune faisant défaut indépendamment avec une probabilité de 5~\% par an (hypothèse) et perdant
80.8~\% des actifs qui y sont détenus (déficit de FTX.com au moment de la requête)~: la perte lorsqu'une plateforme fait défaut, la probabilité d'au moins une défaillance dans l'année, et la
perte annuelle espérée. Données~: \texttt{s1\_crypto.venues}.}\label{fig:s1:crypto-medium-frequency-strategies:venues}
\end{omfigure}

\section{Fiches de stratégie}

\begin{strategyfile}{Carry de financement perpétuel}\label{strat:s1:crypto-medium-frequency-strategies:funding}
\sfield{Qui paie, et pourquoi} Des positions longues à effet de levier qui paient le financement pour détenir des contrats perpétuels.
\sfield{Instruments et plateformes} Cryptomonnaies au comptant et leurs contrats perpétuels sur les plus grandes plateformes.
\sfield{Signal} Taux de financement espéré positif~; sortie lorsqu'il devient durablement négatif.
\sfield{Dimensionnement et exécution} Position longue au comptant, position courte sur le contrat perpétuel à notionnel égal~; marge maintenue au-dessus de la marge de maintenance, avec un coussin.
\sfield{Coûts} Frais de négociation, emprunt si le comptant est financé, et \omterm{def:s1:crypto-medium-frequency-strategies:venue}{risque de plateforme}.
\sfield{Comment elle s'éteint} Le capital d'arbitrage comprime le taux de financement~; défaillance d'une plateforme.
\sfield{Horizon, capacité, infrastructure} Plusieurs mois~; capacité limitée par les positions ouvertes~; suivi de la conservation et de la marge.
\sfield{Backtest honnête} Financement tel que versé, intervalle par intervalle~; frais~; le capital effectivement immobilisé~; une charge pour pertes de plateforme.
\sfield{Sources} Données de financement de Binance (livre~3)~; ce chapitre~: 25.3~\% du capital en 2021, 3.3~\% en 2022, avant \omterm{def:s1:crypto-medium-frequency-strategies:venue}{risque de plateforme}.
\end{strategyfile}

\begin{strategyfile}{Carry de base sur contrats à échéance}\label{strat:s1:crypto-medium-frequency-strategies:basis}
\sfield{Qui paie, et pourquoi} Investisseurs qui veulent une exposition longue à effet de levier via des contrats à terme.
\sfield{Instruments et plateformes} Contrats à échéance sur des plateformes réglementées et des plateformes crypto, contre le comptant.
\sfield{Signal} La base annualisée jusqu'à l'échéance.
\sfield{Dimensionnement et exécution} Position longue au comptant, position courte sur le contrat à échéance, conservée jusqu'à l'échéance.
\sfield{Coûts} Frais~; capital pour la marge~; le roll vers l'échéance suivante.
\sfield{Comment elle s'éteint} Davantage de capital d'arbitrage~; une réglementation qui l'allège.
\sfield{Horizon, capacité, infrastructure} Quelques semaines à quelques mois.
\sfield{Backtest honnête} Base aux prix exécutables~; marge selon les règles de la plateforme.
\sfield{Sources} Schmeling, Schrimpf et Todorov~: carry parfois supérieur à 40~\% par an.
\end{strategyfile}

\begin{strategyfile}{Momentum transversal sur les cryptomonnaies}\label{strat:s1:crypto-medium-frequency-strategies:momentum}
\sfield{Qui paie, et pourquoi} Des investisseurs qui sous-réagissent, puis suivent la tendance~; flux d'attention.
\sfield{Instruments et plateformes} Comptant ou contrats perpétuels sur les cryptomonnaies liquides.
\sfield{Signal} Rendements passés de trois à huit semaines, classés.
\sfield{Dimensionnement et exécution} Quintiles longs--courts hebdomadaires, pondérés par la volatilité.
\sfield{Coûts} Élevés pour les petites cryptomonnaies.
\sfield{Comment elle s'éteint} Effondrements et retournements~; radiations.
\sfield{Horizon, capacité, infrastructure} Semaines~; historiques de prix nettoyés, incluant les cryptomonnaies disparues.
\sfield{Backtest honnête} Biais de survie~: inclure les cryptomonnaies radiées~; coûts par cryptomonnaie.
\sfield{Sources} Liu et Tsyvinski (2021)~; Liu, Tsyvinski et Wu (2022)~; ce chapitre~: 1.35 après 10 pb à quatre semaines, synthétique.
\end{strategyfile}

\begin{strategyfile}{Signal de flux on-chain}\label{strat:s1:crypto-medium-frequency-strategies:onchain}
\sfield{Qui paie, et pourquoi} Détenteurs sur le point de vendre, dont les dépôts vers les bourses sont publics avant leurs ordres.
\sfield{Instruments et plateformes} Cryptomonnaies dont les adresses des bourses sont bien étiquetées.
\sfield{Signal} Entrées nettes vers les bourses, en score z.
\sfield{Dimensionnement et exécution} Positions tactiques à court horizon~; de petite taille.
\sfield{Coûts} Étiquetage des données~; rotation du portefeuille.
\sfield{Comment elle s'éteint} Transferts internes pris pour des dépôts~; usage plus large des mêmes données.
\sfield{Horizon, capacité, infrastructure} Jours~; un nœud ou un fournisseur de données, et des étiquettes d'adresses.
\sfield{Backtest honnête} Étiquettes telles qu'elles étaient connues à l'époque~; horodatages de bloc, et non horodatages indexés a posteriori.
\sfield{Sources} Aucun chiffre de performance vérifié~; le test synthétique du chapitre.
\end{strategyfile}

\begin{strategyfile}{Retour à la moyenne du taux de financement}\label{strat:s1:crypto-medium-frequency-strategies:meanrev}
\sfield{Qui paie, et pourquoi} Positions à effet de levier encombrées qui se débouclent.
\sfield{Instruments et plateformes} Contrats perpétuels.
\sfield{Signal} Taux de financement très au-dessus ou très en dessous de son intervalle récent.
\sfield{Dimensionnement et exécution} Parier contre les taux de financement extrêmes, de taille réduite.
\sfield{Coûts} Financement versé pendant l'attente~; cascades de liquidations.
\sfield{Comment elle s'éteint} Tendances qui persistent tant que le financement reste extrême.
\sfield{Horizon, capacité, infrastructure} Jours.
\sfield{Backtest honnête} Financement à chaque intervalle de financement~; risque de liquidation selon les règles de marge de la plateforme.
\sfield{Sources} Aucun chiffre de performance vérifié.
\end{strategyfile}

\begin{strategyfile}{Carry diversifié entre plateformes}\label{strat:s1:crypto-medium-frequency-strategies:venue}
\sfield{Qui paie, et pourquoi} Comme pour le \omterm{def:s1:crypto-medium-frequency-strategies:carry}{carry de financement}.
\sfield{Instruments et plateformes} La même opération sur plusieurs plateformes.
\sfield{Signal} Taux de financement par plateforme, net du risque propre à chacune.
\sfield{Dimensionnement et exécution} Capital réparti selon la qualité des plateformes~; plafonds de retrait et seuils de déclenchement.
\sfield{Coûts} Moins de compensation des positions, davantage de capital inactif.
\sfield{Comment elle s'éteint} Défaillances corrélées lorsque l'effondrement d'une plateforme se propage.
\sfield{Horizon, capacité, infrastructure} Mois~; opérations de trésorerie entre plateformes.
\sfield{Backtest honnête} Une charge pour défaillance de plateforme~; défaillances présentes dans l'échantillon.
\sfield{Sources} Analyse des débiteurs de FTX.com (Doc 792-1)~; ce chapitre~: 20.2~\% du capital perdu lorsqu'une des quatre plateformes fait défaut.
\end{strategyfile}

\section{Tutoriel~: jusqu'à la défaillance de la bourse}\label{tut:s1:crypto-medium-frequency-strategies:tut}

\begin{tutorial}
\textbf{Objectif.} Transformer les statistiques de financement de Binance en rendements annuels d'un portefeuille de carry, mesurer la perte de FTX en cas de défaillance à partir de l'analyse de ses débiteurs,
simuler le momentum des jetons et les flux vers les bourses, et dimensionner un portefeuille de carry entre plateformes. \textbf{État final~:} les tableaux et les deux figures.
\begin{tutsteps}
\item \textbf{Les données réelles}~: financement par année, soldes et flux de FTX.
\omcode{code/strategies-1/27-crypto-medium-frequency-strategies/python/s1_crypto.py}{38}{54}{Carry de financement par année et déficit de FTX.}\label{lst:s1:crypto-medium-frequency-strategies:data}
\item \textbf{Momentum, flux et plateformes}.
\omcode{code/firm/cryptomf/firm_cryptomf.py}{65}{100}{Momentum des jetons, IC des flux et pertes par plateforme.}\label{lst:s1:crypto-medium-frequency-strategies:venues}
\item \textbf{Exécuter} \texttt{funding()}, \texttt{ftx()}, \texttt{momentum\_table()}, \texttt{flows()} et \texttt{venues()}, ainsi que \texttt{fig\_crypto.py}.
\end{tutsteps}
\textbf{Pour aller plus loin.} Corréler les défaillances de plateformes (une défaillance augmente la probabilité des autres)~; récupérer les taux de financement intervalle par intervalle et calculer
les drawdowns~; ajouter un processus de radiation à l'univers de jetons et observer le biais de survie du momentum.
\end{tutorial}

\section{Construction~: stratégies crypto}\label{bld:s1:crypto-medium-frequency-strategies:cryptomf}

\begin{build}
\bfield{Rôle} Comptabilité du carry, un univers de jetons avec momentum et flux, et pertes liées à la défaillance d'une plateforme.
\bfield{Interface} \texttt{carry\_return(funding, margin, cost)}, \texttt{basis\_carry(fut, spot, days)}, \texttt{CoinConfig(\dots)},
\texttt{simulate\_coins(cfg, rng)}, \texttt{momentum\_book(r, lookback, hold, q, cost)}, \texttt{flow\_ic(flow, r, window)},
\texttt{venue\_losses(k, p, lgd, years, rng)}.
\bfield{Règles} Carry sur le capital immobilisé~; signaux fondés sur des données passées~; défaillances de plateformes indépendantes sauf mention contraire.
\bfield{Tests d'acceptation} \texttt{code/firm/cryptomf/tests/}~: carry et base à la main, momentum et IC des flux sur un univers artificiel, la moyenne et le support
des pertes par plateforme.
\bfield{Pour aller plus loin} Défaillances corrélées~; financement par intervalle~; radiations.
\end{build}

\begin{omsources}
\item Y.~Liu et A.~Tsyvinski, «~Risks and returns of cryptocurrency~», \emph{Review of Financial Studies} 34(6), 2021.
\item Y.~Liu, A.~Tsyvinski et X.~Wu, «~Common risk factors in cryptocurrency~», \emph{Journal of Finance} 77(2), 2022.
\item M.~Schmeling, A.~Schrimpf et K.~Todorov, «~Crypto carry~», \emph{Management Science}, 2026.
\item FTX Debtors, \emph{Preliminary Analysis of Shortfalls}, Doc 792-1, In re FTX Trading Ltd., 22-11068 (Bankr. D. Del.), mars 2023.
\item Binance, taux de financement du contrat perpétuel BTCUSDT (point d'accès public) et documentation du taux de financement.
\end{omsources}

\section{Exercices}

\begin{exercise}[$\star$]\label{exo:s1:crypto-medium-frequency-strategies:1}
Donner l'équivalent annualisé d'un taux de financement de 0.01~\% par intervalle de huit heures, versé trois fois par jour, puis celui de 0.1~\%.
\end{exercise}

\begin{exercise}[$\star$]\label{exo:s1:crypto-medium-frequency-strategies:2}
Le financement moyen s'élève à 30.61~\% par an~; les coûts sont de 0.2~\% par an~; la marge du contrat perpétuel représente 20~\% du notionnel. Calculer le rendement sur capital du portefeuille.
\end{exercise}

\begin{exercise}[$\star$]\label{exo:s1:crypto-medium-frequency-strategies:3}
Expliquer pourquoi un portefeuille acheteur au comptant et vendeur à découvert du contrat perpétuel n'est pas exposé au cours du jeton, et à quoi il est exposé.
\end{exercise}

\begin{exercise}[$\star\star$]\label{exo:s1:crypto-medium-frequency-strategies:4}
Les plateformes font défaut indépendamment, avec une probabilité de 5~\% par an, et perdent 80.8~\% des actifs qui s'y trouvent. Calculer la perte annuelle espérée et la probabilité qu'au moins
une défaillance parmi quatre plateformes, puis parmi huit.
\end{exercise}

\begin{exercise}[$\star\star$]\label{exo:s1:crypto-medium-frequency-strategies:5}
Calculer le rendement du portefeuille de carry en 2021 et en 2022, net de la perte espérée liée aux plateformes, fixée à 4.0~\% par an.
\end{exercise}

\begin{exercise}[$\star\star$]\label{exo:s1:crypto-medium-frequency-strategies:6}
Expliquer pourquoi le portefeuille de momentum à une semaine perd au-delà de 20 points de base, alors que celui à quatre semaines ne perd pas.
\end{exercise}

\begin{exercise}[$\star\star\star$]\label{exo:s1:crypto-medium-frequency-strategies:7}
\emph{Programmation.} Exécuter \texttt{venues\_correlated()}~: une année où une plateforme fait défaut, chaque autre plateforme fait aussi défaut avec une probabilité de 30~\%. Comparer la perte espérée
et le centile de rang 99 avec le cas indépendant, pour quatre plateformes.
\end{exercise}

\begin{exercise}[$\star\star\star$]\label{exo:s1:crypto-medium-frequency-strategies:8}
\emph{Trouver la faille.} «~Notre \omterm{def:s1:crypto-medium-frequency-strategies:carry}{carry de financement} a rapporté 18~\% par an pendant cinq ans, avec un drawdown maximal de 2~\%~: il est neutre au marché et presque sans risque.~»
\end{exercise}

\section{Problème~: jusqu'à la défaillance de la bourse}

\begin{problem}[{Problème du week-end --- carry, momentum et plateforme}]\label{pb:s1:crypto-medium-frequency-strategies:1}
Les données de financement de Binance, l'analyse des débiteurs de FTX, les jetons synthétiques et le dossier public.

\medskip\noindent\textbf{Partie I --- Carry.}
\begin{enumerate}
\item Définir le \omterm{def:s1:crypto-medium-frequency-strategies:carry}{carry de financement} et le \omterm{def:s1:crypto-medium-frequency-strategies:carry}{carry de base}.
\item Comment le taux de financement de Binance est-il fixé, et à quel moment est-il versé~?
\item Donner le rendement sur capital du portefeuille de carry, année par année.
\item Que montrent Schmeling, Schrimpf et Todorov sur le carry crypto~?
\end{enumerate}

\medskip\noindent\textbf{Partie II --- Momentum et flux.}
\begin{enumerate}[resume]
\item Que montrent Liu et Tsyvinski, puis Liu, Tsyvinski et Wu~?
\item Décrire l'univers synthétique de jetons.
\item Donner les ratios de Sharpe du portefeuille de momentum selon l'horizon de formation et le coût.
\item Définir un \omterm{def:s1:crypto-medium-frequency-strategies:onchain}{signal on-chain}~; que révèle le test sur les flux, et que ne montre-t-il pas~?
\end{enumerate}

\medskip\noindent\textbf{Partie III --- Plateformes.}
\begin{enumerate}[resume]
\item Définir le \omterm{def:s1:crypto-medium-frequency-strategies:venue}{risque de plateforme}.
\item Que montrent les soldes de FTX.com à la date du dépôt de bilan et les flux de novembre~?
\item Donner la perte lorsqu'une des $k$ plateformes fait défaut, et la probabilité d'une défaillance par an.
\item Pourquoi la diversification ne modifie-t-elle pas la perte espérée~?
\end{enumerate}

\medskip\noindent\textbf{Partie IV --- Le verdict.}
\begin{enumerate}[resume]
\item Énoncer le \emph{résultat nommé}~: le rendement annuel du \omterm{def:s1:crypto-medium-frequency-strategies:carry}{carry de financement} et la perte du portefeuille diversifié entre plateformes lorsqu'une plateforme fait défaut.
\item Qu'est-ce qui réduit la perte espérée liée aux plateformes~?
\item Comment backtester honnêtement le \omterm{def:s1:crypto-medium-frequency-strategies:carry}{carry de financement}~?
\item Quelle fiche de stratégie porte le plus de \omterm{def:s1:crypto-medium-frequency-strategies:venue}{risque de plateforme}~?
\item Quel est le lien entre le carry de ce chapitre et celui du chapitre~20~?
\item Que rémunérait le carry en 2021~?
\item Pourquoi le financement a-t-il baissé après 2022~?
\item En une phrase~: pour quoi un trader de carry crypto est-il rémunéré~?
\end{enumerate}
\end{problem}

\section{Questions d'entretien}

\begin{interviewq}[$\star$ \iqroles{chercheur}]\label{iq:s1:crypto-medium-frequency-strategies:1}
Qu'est-ce que le taux de financement d'un contrat perpétuel, et pourquoi existe-t-il~?
\end{interviewq}

\begin{interviewq}[$\star\star$ \iqroles{trader}]\label{iq:s1:crypto-medium-frequency-strategies:2}
Le financement d'un jeton s'élève à 60~\% par an. Décrivez l'opération et ses risques.
\end{interviewq}

\begin{interviewq}[$\star\star$ \iqroles{risque}]\label{iq:s1:crypto-medium-frequency-strategies:3}
Comment fixeriez-vous des limites sur les actifs qu'un établissement détient sur chaque plateforme crypto~?
\end{interviewq}

\begin{interviewq}[$\star\star$ \iqroles{développeur}]\label{iq:s1:crypto-medium-frequency-strategies:4}
Que doit surveiller un système qui fait tourner un portefeuille de carry réparti sur plusieurs plateformes~?
\end{interviewq}

\begin{interviewq}[$\star\star$ \iqroles{chercheur}]\label{iq:s1:crypto-medium-frequency-strategies:5}
Comment testeriez-vous le momentum des jetons sans biais de survie~?
\end{interviewq}

\begin{interviewq}[$\star\star\star$ \iqroles{chercheur}]\label{iq:s1:crypto-medium-frequency-strategies:6}
Un portefeuille répartit son capital sur $k$ plateformes, chacune faisant défaut indépendamment avec une probabilité $p$ par an et perdant une fraction $L$ de ses actifs. Établissez la moyenne et la variance
de la perte annuelle, et la probabilité qu'elle soit au moins égale à $mL/k$ pour un entier $m$.
\end{interviewq}
